\chapter{Related Work}
\label{chap:related_work}

\section{Safety Alignment in Large Language Models}

Modern LLMs are trained to refuse harmful requests through a post-training alignment stage. The dominant approach, Reinforcement Learning from Human Feedback (RLHF)~\citep{ouyang2022instructgpt}, fine-tunes a language model using human preference judgements that reward helpful and harmless responses. Constitutional AI~\citep{bai2022constitutional} extends this with a set of explicit principles that guide self-critique and revision. Both approaches have been adopted across major models, including GPT-4~\citep{openai2023gpt4}, Claude~\citep{anthropic2024claude}, Llama~\citep{llama31}, and Gemma~\citep{gemma2}.

A key empirical observation across all aligned models is that alignment quality is strongly correlated with the language of the interaction. Safety benchmarks are predominantly designed and curated in English~\citep{gehman2020realtoxicityprompts, harmbench}, and the human feedback collected during RLHF is likewise skewed toward English-speaking annotators. As a result, aligned models exhibit substantially stronger refusal behavior in English than in other languages~\citep{deng2023multilingual, yong2023lowresource, wang2023allsafe}, a gap that widens further for languages with limited digital presence.

\section{Jailbreak Attacks}

Jailbreak attacks are inputs designed to bypass a model's safety guardrails and elicit harmful responses. They can be broadly categorised by their construction strategy.

\paragraph{Prompt-based attacks.} Early jailbreaks exploited manually crafted prompts such as roleplay instructions (``pretend you are DAN'') or hypothetical framings to sidestep safety training~\citep{perez2022jailbreak}. These attacks highlighted that aligned models could be manipulated by altering the framing of a request without changing its underlying harmful content.

\paragraph{Optimisation-based attacks.} Zou et al.~\citeyearpar{zou2023gcg} introduced GCG (Greedy Coordinate Gradient), which automatically optimises a short adversarial suffix appended to a harmful prompt to maximise the probability of compliance. AutoDAN~\citep{liu2023autodan} generates jailbreaks using genetic algorithms. These methods demonstrate that sufficiently flexible prompt manipulation can reliably bypass aligned models, though the resulting prompts are often semantically anomalous.

\paragraph{Representation-level attacks.} A more recent line of work directly intervenes on the model's internal representations. Representation engineering~\citep{zou2023representation} showed that adding a direction associated with a target behavior to intermediate activations can steer model outputs. Arditi et al.~\citeyearpar{arditi2024refusal} demonstrated that adding the negation of the refusal direction at a single layer reliably induces compliance with harmful requests, providing a mechanistically grounded attack method.

Our work extends this line of inquiry to the multilingual setting by extracting jailbreak-specific directions $\mathbf{j}_l$ from compliance behavior, rather than from the contrast between harmful and harmless prompts.

\section{Multilingual Jailbreak}
\label{sec:multilingual_jailbreak}

The observation that non-English languages weaken safety guardrails emerged empirically before mechanistic explanations were available. Several parallel lines of work have documented this phenomenon.

\paragraph{Low-resource language vulnerability.}
Yong et al.~\citeyearpar{yong2023lowresource} found that translating harmful English prompts into low-resource languages such as Scots Gaelic, Zulu, or Hmong reliably bypassed GPT-4's content filters. They attribute this to the sparse representation of these languages in alignment training data. Shen et al.~\citeyearpar{shen2024language} demonstrated similar effects across a broader range of models, finding that safety alignment degrades consistently as a function of language resource level. These findings suggest that LRL vulnerability is not model-specific but reflects a systemic property of how alignment is distributed across languages.

\paragraph{Multilingual attack diversity.}
Deng et al.~\citeyearpar{deng2023multilingual} conducted a systematic evaluation across 54 languages, showing that multilingual jailbreak success rates are inversely correlated with a language's representation in model training data. They also found that combining languages — for example, mixing English and Chinese in a single prompt — can further reduce refusal rates, suggesting that language-switching disrupts alignment mechanisms rather than simply exploiting coverage gaps.

\paragraph{Multilingual safety benchmarks.}
Several benchmarks have been proposed to measure safety across languages. SafetyBench-ML~\citep{zhang2023safetybench} covers ten languages but is limited to multiple-choice evaluation. PolyRefuse~\citep{polyrefuse} provides full generation evaluation across 14 languages with WildGuard scoring, enabling measurement of refusal rates as a continuous signal. Our work extends PolyRefuse to 16 languages by adding Swahili and Amharic, and uses it as the primary evaluation dataset.

\paragraph{Mechanistic accounts.}
While the empirical picture of multilingual jailbreak is well-established, mechanistic explanations remain limited. The most directly relevant prior work is that of Arditi et al.~\citeyearpar{arditi2024refusal}, who showed that the refusal direction $\mathbf{r}_l$ is geometrically consistent across safety-aligned HRL languages and that subtracting it at inference time reliably induces compliance — effectively a same-language jailbreak attack. Their appendix (Table~6 and Table~7) demonstrates the feasibility of same-language $\pm \mathbf{r}_l$ interventions but does not study cross-lingual transfer, does not analyse LRL languages, and does not define a jailbreak-specific direction. Our work addresses all three of these dimensions.

\section{Representation Engineering and Mechanistic Interpretability}

Mechanistic interpretability studies how model behaviors are encoded in internal representations. A key finding is that many semantic properties — from factual associations to emotional content — are linearly encoded as directions in the residual stream~\citep{elhage2022superposition, burns2022discovering}. This motivates the use of linear probes and direction-based interventions as a tool for studying and controlling model behavior.

The Difference-in-Means (DiM) method, introduced by~\citet{zou2023representation} as part of the Representation Engineering framework, computes a direction that separates two contrastive sets of activations (e.g., harmful vs.\ harmless inputs). DiM directions have been used to detect lying behavior~\citep{marks2023geometry}, suppress refusal~\citep{arditi2024refusal}, and induce target emotions~\citep{zou2023representation}. Our jailbreak vector $\mathbf{j}_l$ applies DiM to a different contrast: bypassed vs.\ refused completions of the same harmful inputs, extracted entirely from baseline inference.

\section{Cross-lingual Transfer in LLMs}

LLMs trained on multilingual corpora develop shared internal representations across languages, enabling a degree of cross-lingual transfer. Pires et al.~\citeyearpar{pires2019multilingual} showed that multilingual BERT encodes semantically similar sentences in nearby regions of activation space regardless of language. Subsequent work confirmed that this cross-lingual alignment extends to instruction-tuned models, where knowledge and reasoning abilities transfer across languages~\citep{shi2022language, zhao2024llm}.

In the safety domain, cross-lingual transfer has been studied primarily at the surface level — for example, testing whether a model fine-tuned on English safety data generalises to Chinese refusal~\citep{wang2023allsafe}. At the representation level, it remains open whether safety vectors transfer across languages. The base paper~\citep{arditi2024refusal} provides indirect evidence by showing that $\mathbf{r}_l$ is similar across HRL languages, but does not test whether $\mathbf{r}_l$ from one language can steer behavior in another. Our experiments directly test this for the jailbreak direction $\mathbf{j}_l$, providing the first explicit cross-lingual transfer results at the activation level.

\section{Summary and Positioning}

Table~\ref{tab:related_work} positions this work relative to the most directly relevant prior studies.

\begin{table}[h]
\centering
\caption{Comparison with related work. \checkmark = addressed; $\circ$ = partially addressed; -- = not addressed.}
\label{tab:related_work}
\setlength{\tabcolsep}{5pt}
\begin{tabular}{lcccccc}
\toprule
\textbf{Work} & \textbf{Multilingual} & \textbf{LRL} & \textbf{Jailbreak vec.} & \textbf{Mechanism} & \textbf{Cross-lingual} & \textbf{Defense} \\
\midrule
Arditi et al.~\citeyearpar{arditi2024refusal}  & $\circ$   & --         & --         & \checkmark & --         & \checkmark \\
Yong et al.~\citeyearpar{yong2023lowresource}   & \checkmark & \checkmark & --         & --         & --         & --         \\
Deng et al.~\citeyearpar{deng2023multilingual}  & \checkmark & \checkmark & --         & --         & --         & --         \\
Zou et al.~\citeyearpar{zou2023representation}  & --         & --         & $\circ$    & \checkmark & --         & --         \\
\midrule
\textbf{This work}                               & \checkmark & \checkmark & \checkmark & \checkmark & \checkmark & \checkmark \\
\bottomrule
\end{tabular}
\end{table}

The key distinction of our approach is that it operates at the intersection of three dimensions that have previously been studied in isolation: mechanistic interpretability of safety representations, multilingual coverage including LRL languages, and cross-lingual transfer of both attack and defense vectors. By extracting jailbreak vectors purely from inference behavior — without requiring harmless data or training — we also provide a method that is applicable in the data-sparse settings characteristic of low-resource languages.
